% TEMPLATE-MILC-v3.tex - plantilla maestra para todas las guias MILC v3
% Ver scripts/generadores/build-guia-milc.py para la lista completa de
% claves esperadas. El script valida que no queden placeholders sin
% reemplazar antes de escribir el archivo final.

\documentclass[11pt,letterpaper]{article}
\usepackage[margin=1.75cm]{geometry}
\usepackage[table]{xcolor}
\usepackage{fontspec}
\usepackage[spanish]{babel}
\usepackage{tikz}
\usepackage[most]{tcolorbox}
\usepackage{tabularx}
\usepackage{array}
\usepackage{fancyhdr}
\usepackage{titlesec}
\usepackage{ragged2e}
\usepackage{enumitem}
\usepackage{microtype}

\setmainfont{Helvetica Neue}
\setsansfont{Helvetica Neue}
\setlength{\parindent}{0pt}
\setlength{\parskip}{6pt}
\hyphenpenalty=200
\exhyphenpenalty=200
\pretolerance=400
\tolerance=2000
\setlength{\emergencystretch}{1em}
\renewcommand{\arraystretch}{1.25}
\setlist{leftmargin=5mm,itemsep=1mm,topsep=1mm}
\newcolumntype{Y}{>{\RaggedRight\arraybackslash}X}

\definecolor{milcMagenta}{HTML}{E6007E}
\definecolor{milcVino}{HTML}{5A0038}
\definecolor{milcTurquesa}{HTML}{0093A5}
\definecolor{milcVerde}{HTML}{58A923}
\definecolor{milcMostaza}{HTML}{E5B400}
\definecolor{milcOcre}{HTML}{B86600}
\definecolor{milcNegro}{HTML}{241816}
\definecolor{milcGris}{HTML}{F7F7F4}
\definecolor{milcLinea}{HTML}{E8E2DD}
\definecolor{milcGradePrimary}{HTML}{84CC16}
\definecolor{milcGradeDark}{HTML}{4D7C0F}
\definecolor{milcGradeSoft}{HTML}{ECFCCB}
\definecolor{milcGradeAccent}{HTML}{CFE7FF}
\definecolor{milcGradeText}{HTML}{FFFFFF}
\color{milcNegro}

\pagestyle{fancy}
\fancyhf{}
\lhead{\small Tecnología e Informática · Grado 7}
\rhead{\small Guía 29 · MILC v3}
\cfoot{\small\thepage}
\renewcommand{\headrulewidth}{0pt}

\newtcolorbox{titlebox}[1]{
  enhanced,colback=#1,colframe=#1,arc=7mm,boxrule=0pt,
  left=7mm,right=7mm,top=3mm,bottom=3mm,
  before skip=8pt,after skip=8pt,
  fontupper=\sffamily\bfseries\Large\color{white}
}

\newtcolorbox{softbox}[3]{
  enhanced,colback=#2,colframe=#1,borderline west={4pt}{0pt}{#1},
  arc=2mm,boxrule=.4pt,left=4mm,right=4mm,top=3mm,bottom=3mm,
  before skip=6pt,after skip=8pt,title={#3},coltitle=white,
  colbacktitle=#1,fonttitle=\sffamily\bfseries,fontupper=\RaggedRight,
  attach boxed title to top left={xshift=3mm,yshift=-2mm},
  boxed title style={arc=2mm, boxrule=0pt}
}

\newtcolorbox{drawbox}[2]{
  enhanced,colback=white,colframe=#1,arc=4mm,boxrule=1pt,
  left=5mm,right=5mm,top=4mm,bottom=4mm,before skip=8pt,after skip=8pt,
  title={#2},coltitle=white,colbacktitle=#1,fonttitle=\sffamily\bfseries,
  fontupper=\RaggedRight,
  attach boxed title to top left={xshift=4mm,yshift=-2mm},
  boxed title style={arc=3mm, boxrule=0pt}
}

\newtcolorbox{infoband}[1]{
  enhanced,colback=#1!8,colframe=#1!50,arc=2mm,boxrule=.5pt,
  left=4mm,right=4mm,top=2mm,bottom=2mm,before skip=4pt,after skip=4pt,
  fontupper=\small\RaggedRight
}

% Puente narrativo entre secciones (contrato editorial Conéctate).
% Da continuidad: "Acabas de X. Ahora vas a Y porque Z."
% Visualmente: banda izquierda en color del grado, texto en itálica pequeña.
\newtcolorbox{puentebox}{
  enhanced,colback=milcGradeSoft,colframe=milcGradeDark,
  borderline west={3pt}{0pt}{milcGradeDark},
  arc=0mm,boxrule=0pt,left=5mm,right=4mm,top=2.5mm,bottom=2.5mm,
  before skip=4pt,after skip=8pt,
  fontupper=\itshape\small\color{milcNegro}\RaggedRight
}
\newcommand{\puente}[1]{%
  \begin{puentebox}%
    \textcolor{milcGradeDark}{\textbf{\textrightarrow}}\hspace{2mm}#1%
  \end{puentebox}%
}

\newcommand{\linead}[1]{\noindent\color{milcLinea}\rule{#1}{0.5pt}\color{milcNegro}}
\newcommand{\respuesta}[1]{\par\vspace{2mm}\foreach \n in {1,...,#1}{\noindent\color{milcLinea}\rule{\linewidth}{0.5pt}\par\vspace{5mm}}\color{milcNegro}}
\newcommand{\checkbox}{\textcolor{milcGradeDark}{\fbox{\phantom{x}}}\hspace{5pt}}

\newcommand{\phasecard}[4]{
\begin{tcolorbox}[enhanced,colback=#3,colframe=#2,arc=3mm,boxrule=.5pt,height=3.05cm,valign=center,halign=center,left=1.5mm,right=1.5mm,top=2mm,bottom=2mm]
{\sffamily\bfseries\fontsize{22}{24}\selectfont\textcolor{#2}{#1}}\par
{\sffamily\bfseries\small #4}
\end{tcolorbox}
}

\begin{document}
\thispagestyle{empty}
\begin{tikzpicture}[remember picture,overlay]
  \fill[white] (current page.south west) rectangle (current page.north east);
  \path[fill=milcGradePrimary,rounded corners=16mm]
    ([xshift=.55cm,yshift=-.55cm]current page.north west) rectangle
    ([xshift=-.55cm,yshift=3.65cm]current page.south east);

  \node[anchor=north west,text=milcGradeText] at ([xshift=2.0cm,yshift=-1.85cm]current page.north west)
    {\sffamily\bfseries\fontsize{14}{16}\selectfont GRADO};
  \node[anchor=north west,text=milcGradeText,opacity=.82] at ([xshift=2.0cm,yshift=-2.75cm]current page.north west)
    {\sffamily\bfseries\fontsize{9}{11}\selectfont SERIE GUÍAS · TECNOLOGÍA E INFORMÁTICA};

  \begin{scope}[shift={([xshift=-3.05cm,yshift=-2.55cm]current page.north east)}]
    \fill[white,opacity=.80] (-.62,-.52) rectangle (.62,.52);
    \draw[milcGradeText,line width=1pt] (-.62,-.52) rectangle (.62,.52);
    \fill[milcGradeDark] (-.42,-.38) rectangle (-.22,.06);
    \fill[milcGradeAccent] (-.10,-.38) rectangle (.10,.24);
    \fill[milcGradePrimary!60!black] (.22,-.38) rectangle (.42,.38);
  \end{scope}

  \node[anchor=west,text=milcGradeText] at ([xshift=1.9cm,yshift=-12.85cm]current page.north west)
    {\sffamily\bfseries\fontsize{124}{124}\selectfont 7};
  \draw[milcGradeText,line width=1.7pt]
    ([xshift=4.52cm,yshift=-11.68cm]current page.north west) circle (.13cm);

  \node[anchor=west,text=milcGradeText,text width=8.7cm,align=left] at ([xshift=10.35cm,yshift=-11.6cm]current page.north west)
    {
     {\sffamily\bfseries\fontsize{19}{22}\selectfont IA y\\mis tareas ---\\política\\personal\\responsable}\\[4mm]
     {\sffamily\bfseries\fontsize{23}{26}\selectfont Guía 29-7-TIC}\\[2mm]
     {\sffamily\fontsize{13}{16}\selectfont Período 3 · Inteligencia Artificial}\\[5mm]
     {\sffamily\bfseries\fontsize{11}{13}\selectfont Institución Educativa Sor María Juliana}\\[1mm]
     {\sffamily\fontsize{10}{12}\selectfont Autor: Dr. Álvaro Cárdenas Orozco}
    };

  \path[fill=milcGradeSoft,rounded corners=7mm]
    ([xshift=1.35cm,yshift=.85cm]current page.south west) rectangle
    ([xshift=-1.35cm,yshift=4.25cm]current page.south east);
  \draw[milcGradeDark,line width=.65pt,rounded corners=7mm]
    ([xshift=1.35cm,yshift=.85cm]current page.south west) rectangle
    ([xshift=-1.35cm,yshift=4.25cm]current page.south east);
  \node[anchor=center,text=milcNegro,text width=17.0cm,align=center] at ([yshift=2.55cm]current page.south)
    {\sffamily\fontsize{10}{12}\selectfont Metodología MILC · Apertura ancestral · Pensamiento computacional · Triángulo Dussel-Estoicismo-Floridi\\
    \textcolor{milcGradeDark}{\bfseries Producto:} Tu \textbf{política personal de uso de IA en estudios} (1 hoja del cuaderno): \textbf{(a) los 4 usos LEGÍTIMOS de IA en tareas} con ejemplos concretos; \textbf{(b) los 4 usos NO LEGÍTIMOS (trampa)} con ejemplos; \textbf{(c) cómo atribuir IA correctamente} en un trabajo; \textbf{(d) 5 reglas personales} para tu uso de IA durante el año académico; \textbf{(e) ejemplo aplicado}: un caso real donde usaste IA en tarea (cómo lo manejaste). Firma.};
\end{tikzpicture}
\mbox{}
\newpage

\section*{Apertura: IA y mis tareas --- política personal de uso responsable en estudios}

\begin{softbox}{milcGradeDark}{milcGradeSoft}{Saber ancestral}
\textbf{Doña Mercedes la maestra rural de Cartago decía: ``El que aprende con esfuerzo se gradúa; el que copia, se gradúa también, pero sin oficio.''} En la escuelita de la vereda La Plata había dos niños: \emph{Pedro y Lucía}. Pedro hacía los ejercicios solo, se equivocaba, los corregía, aprendía. Lucía copiaba la tarea de Pedro y la entregaba como propia. \textbf{Las dos pasaron de grado}: las notas no detectaban la diferencia. Pero al año siguiente, cuando llegaron nuevos contenidos que se construían sobre lo anterior, \emph{Pedro entendía; Lucía no}. \emph{Lucía había aprobado, pero no había aprendido}. Tres años después, en un examen de ingreso al colegio de Buga, Pedro pasó; Lucía no. \textbf{La diferencia no estaba en la inteligencia: estaba en haber aprendido el oficio en su momento}. Con IA pasa exactamente lo mismo. \emph{El que la usa para hacer trampa pasa de grado pero no aprende. El que la usa con criterio, asistente real, aprende Y avanza}. Hoy defines tu política personal para no ser Lucía.
\end{softbox}

\begin{softbox}{milcTurquesa}{milcGris}{Saber contemporáneo}
\textbf{Usar IA en tareas escolares es ambiguo}: a veces es legítimo, a veces es trampa. La regla universal: \emph{si TÚ aprendiste, es legítimo; si solo entregaste un producto sin entender, es trampa}. \textbf{Los 4 usos LEGÍTIMOS de IA en tareas:} \textbf{(1) Como tutor:} \emph{``explícame este concepto''}, después tú lo escribes con tus palabras. \textbf{(2) Como corrector:} \emph{``revisa esta carta y dime qué mejorar''}, después tú aplicas los cambios con tu criterio. \textbf{(3) Como organizador:} \emph{``ayúdame a estructurar las ideas para un ensayo sobre X''}, después tú escribes el ensayo. \textbf{(4) Como generador de ideas:} \emph{``dame 5 ángulos posibles para investigar X''}, después tú escoges 1 y profundizas. \textbf{Los 4 usos NO LEGÍTIMOS (trampa):} \textbf{(5) Escribir el ensayo completo} y entregarlo como tuyo. \textbf{(6) Resolver el problema y copiar la respuesta} sin entenderla. \textbf{(7) Hacer toda la investigación} y solo entregar lo que la IA dice, sin verificar ni añadir tu voz. \textbf{(8) Pasar por tu autoría algo que es 100\% generado por IA}. \textbf{La diferencia siempre: ¿aprendí? ¿agregué mi voz?}.
\end{softbox}

\begin{softbox}{milcMostaza}{milcGris}{Pregunta-puente}
\textit{Si tu profe te pide un ensayo de 2 páginas sobre la independencia de Colombia, y tú le pides a ChatGPT que lo escriba completo y se lo entregas tal cual, ¿\textbf{en qué te perjudica}? ¿La nota cambia? ¿Algo más cambia?}
\end{softbox}

\begin{softbox}{milcVerde}{milcGris}{Saber y saber hacer}
Al terminar la clase: (1) podrás \textbf{identificar} los 4 usos legítimos vs los 4 no legítimos; (2) sabrás \textbf{aplicar} la regla universal (¿aprendí?); (3) podrás \textbf{atribuir} IA correctamente en un trabajo; (4) habrás \textbf{creado} tu política personal de uso de IA en estudios.
\end{softbox}

\small
\begin{tabularx}{\textwidth}{>{\bfseries}p{3.0cm}Y}
\rowcolor{milcGradePrimary!12}Autor & Dr. Álvaro Cárdenas Orozco\\
DBA & Reconoce los fundamentos, usos y límites éticos de la Inteligencia Artificial (MEN, T\&I 7°)\\
Referentes & Saber ancestral del consejero · Modelos de lenguaje · Ética algorítmica y prompting\\
Producto final & Tu \textbf{política personal de uso de IA en estudios} (1 hoja del cuaderno): \textbf{(a) los 4 usos LEGÍTIMOS de IA en tareas} con ejemplos concretos; \textbf{(b) los 4 usos NO LEGÍTIMOS (trampa)} con ejemplos; \textbf{(c) cómo atribuir IA correctamente} en un trabajo; \textbf{(d) 5 reglas personales} para tu uso de IA durante el año académico; \textbf{(e) ejemplo aplicado}: un caso real donde usaste IA en tarea (cómo lo manejaste). Firma.\\
\end{tabularx}
\normalsize

\puente{Hoy haces 4 cosas: distingues uso legítimo de trampa, aprendes a atribuir, defines tu política personal, aplicas a un caso real.}

\begin{titlebox}{milcGradeDark}
Ruta MILC: así vamos a aprender
\end{titlebox}
\begin{tabularx}{\textwidth}{>{\centering\arraybackslash}X>{\centering\arraybackslash}X>{\centering\arraybackslash}X>{\centering\arraybackslash}X}
\phasecard{1}{milcVerde}{milcGris}{Escuta\\[-1mm]\small Observo, escucho y pregunto.} &
\phasecard{2}{milcTurquesa}{milcGris}{Sistematización\\[-1mm]\small Pensamiento computacional.} &
\phasecard{3}{milcMagenta}{milcGris}{Praxis\\[-1mm]\small Construyo y aplico.} &
\phasecard{4}{milcVino}{milcGris}{Evaluación\\[-1mm]\small Reflexión liberadora.}
\end{tabularx}


\puente{Empezamos con 8 escenarios para clasificar.}

\begin{titlebox}{milcVerde}
Fase 1 · Escuta
\end{titlebox}

\begin{softbox}{milcVerde}{milcGris}{Escena para escuchar}
\textbf{Actividad 1 · ANALIZA --- 8 escenarios: ¿legítimo o trampa?} (12 min · individual). Lee estos 8 escenarios y marca en el cuaderno LEG (legítimo) o T (trampa), con 1 línea de justificación. \emph{(1) Le pido a ChatGPT que me explique fotosíntesis, después escribo el resumen con mis palabras.} \emph{(2) Le pido a ChatGPT que escriba un ensayo de 2 páginas sobre Bolívar y lo entrego sin cambios.} \emph{(3) Le pido a ChatGPT que me sugiera 5 ángulos para investigar sobre el cambio climático, escojo 1 y desarrollo solo.} \emph{(4) Le pido a ChatGPT que resuelva un problema de matemáticas y copio la respuesta sin entender los pasos.} \emph{(5) Uso ChatGPT para corregir ortografía y gramática de un texto que YO escribí.} \emph{(6) Le pido a ChatGPT que me dé 10 hechos sobre Newton y los pego sin verificar.} \emph{(7) Le pido a ChatGPT que estructure las ideas de mi ensayo, escribo el ensayo aplicando esa estructura.} \emph{(8) Uso ChatGPT para traducir mi tarea de español al inglés, sin saber inglés, y la entrego como prueba de mi nivel de inglés.}
\end{softbox}

\checkbox Leí los 8 escenarios con calma.\\
\checkbox Marqué LEG o T en cada uno.\\
\checkbox {'Identifiqué la regla común': '¿el estudiante aprendió o no?'}

\begin{infoband}{milcVerde}
\textbf{Lo que va al cuaderno (Actividad 1 · ANALIZA).} Título: «Actividad 1 --- 8 escenarios de IA en tareas». \textbf{Formato:} lista numerada con LEG/T + 1 línea de justificación. \textbf{Extensión:} 8 líneas. \textbf{Sabes que terminaste cuando} aplicas la regla universal: ¿aprendí o solo entregué?.
\end{infoband}

\puente{Después aprendes los 4 usos legítimos + los 4 no legítimos + cómo atribuir.}

\begin{titlebox}{milcTurquesa}
Fase 2 · Sistematización con pensamiento computacional
\end{titlebox}

\begin{softbox}{milcTurquesa}{milcGris}{Cuatro pilares aplicados al tema}
Tu intuición se afina ahora. Respuestas correctas: \emph{1, 3, 5, 7 = LEGÍTIMOS. 2, 4, 6, 8 = TRAMPA}. Ahora aprende a fondo los 4 usos legítimos + los 4 no + cómo atribuir.
\end{softbox}

\small
\begin{tabularx}{\textwidth}{>{\bfseries\RaggedRight\arraybackslash}p{3.6cm}Y}
\rowcolor{milcTurquesa!90}\color{white}Pilar & \color{white}\bfseries Aplicación al tema\\
1. Descomposición & \textbf{Los 4 usos LEGÍTIMOS de IA en tareas.} \textbf{(1) Como tutor (explicar):} \emph{``Explícame X''}, después escribo el resumen con mis palabras. La IA ayuda a ENTENDER, no a entregar. \textbf{(2) Como corrector:} \emph{``Revisa este texto que YO escribí y dime qué mejorar''}. La IA detecta errores; tú decides cuáles arreglar. \textbf{(3) Como organizador:} \emph{``Ayúdame a estructurar ideas para un ensayo sobre X''}, después YO escribo el ensayo. La IA ayuda con el ARMAZÓN; el contenido es mío. \textbf{(4) Como generador de ideas (brainstorming):} \emph{``Dame 5 ángulos para investigar X''}, escojo 1 y profundizo solo. La IA abre OPCIONES; yo escojo y desarrollo.\\
2. Reconocimiento de patrones & \textbf{Los 4 usos NO LEGÍTIMOS (trampa).} \textbf{(5) Escribir el ensayo completo} y entregarlo como tuyo. Quitas el aprendizaje, faltas a la honestidad académica. \textbf{(6) Resolver el problema} y copiar respuesta sin entender pasos. No aprendes la metodología. \textbf{(7) Hacer toda la investigación} y solo entregar lo que la IA dice sin verificar ni añadir tu voz. Te haces dependiente. \textbf{(8) Pasar por tu autoría algo 100\% generado por IA.} Es plagio nuevo.\\
3. Abstracción & \textbf{Cómo atribuir IA correctamente en un trabajo.} La práctica profesional emergente es \textbf{citar la IA como herramienta usada}, no como autor. Ejemplos: \emph{(a) En la introducción del trabajo:} \emph{``En la preparación de este ensayo usé Claude AI para estructurar mis ideas iniciales. Las ideas y el desarrollo son míos''}. \emph{(b) Al final del trabajo (sección Fuentes):} \emph{``Asistencia de IA: ChatGPT (versión gratuita), consultado el [fecha] para [propósito]''}. \emph{(c) En tareas escolares específicas:} \emph{``Nota al profe: usé ChatGPT para revisar mi ortografía. El contenido es propio''}. \textbf{La atribución NO es signo de debilidad: es signo de honestidad académica}. Los profesores valoran que digas la verdad.\\
4. Algoritmo conceptual & \textbf{Las 5 preguntas para autoevaluar antes de entregar.} Antes de entregar un trabajo donde usaste IA: \textbf{(1) ¿Aprendí algo de este proceso?} Si no, hay problema. \textbf{(2) ¿Puedo explicar este trabajo si el profe me pregunta verbalmente?} Si no, hay problema. \textbf{(3) ¿Mi voz aparece en el trabajo?} Si solo es voz de IA, hay problema. \textbf{(4) ¿Atribuyé el uso de IA?} Si no, hay problema. \textbf{(5) ¿Mi yo de adulto se sentiría orgulloso de cómo lo hice?} Si no, hay problema. \emph{Si las 5 respuestas son SÍ, el uso fue legítimo}.\\
\end{tabularx}
\normalsize

\begin{drawbox}{milcTurquesa}{4 legítimos + 4 trampa + atribución + 5 preguntas}
\textbf{4 LEGÍTIMOS:} tutor · corrector · organizador · generador de ideas. \\[1mm]
\textbf{4 TRAMPA:} escribir entero · copiar respuesta sin entender · investigar sin verificar/voz · pasar como tuyo lo que es 100\% IA. \\[1mm]
\textbf{Atribuir:} ``En este trabajo usé [IA] para [propósito específico]''. \\[1mm]
\textbf{5 preguntas autoevaluación:} ¿aprendí? ¿puedo explicar? ¿mi voz está? ¿atribuí? ¿yo adulto orgulloso?
\end{drawbox}

\begin{softbox}{milcMostaza}{milcGris}{Errores comunes y su corrección}
\textbf{Errores frecuentes al usar IA en tareas.} (1) \emph{No saber distinguir uso legítimo de trampa}: caes en trampa sin saber. (2) \emph{No atribuir cuando usaste IA}: parece deshonesto aunque el uso haya sido legítimo. (3) \emph{Esperar a estar en problemas para definir tu política}: la política se define en frío, antes. (4) \emph{Asumir que el profe sabe distinguir IA de tu escritura}: cada vez más detectan; mejor ser honesto. (5) \emph{Justificar trampa como ``todos lo hacen''}: que otros la hagan no la hace legítima.
\end{softbox}

\begin{infoband}{milcTurquesa}
\textbf{Lo que va al cuaderno (Actividad 2 · EVALÚA).} Título: «Actividad 2 --- Diagnóstico de mis usos pasados de IA». \textbf{Formato:} bloque A con 4 legítimos (cada uno: descripción + ejemplo). Bloque B con 4 no legítimos. Bloque C con plantilla de atribución. Bloque D con las 5 preguntas de autoevaluación. \textbf{Veredicto honesto (3 líneas):} piensa en las últimas 3 veces que usaste IA para una tarea escolar. Pasa cada uso por las 5 preguntas. ¿Cuántas fueron \emph{legítimas} y cuántas cruzaron la línea? Reconoce sin justificar. \textbf{Extensión:} 18 líneas + veredicto honesto. \textbf{Sabes que terminaste cuando} pudiste decir ``en una de esas crucé la línea'' sin justificar --- ese reconocimiento es el inicio del oficio honesto.
\end{infoband}

\puente{Con todo claro, escribes tu política personal de IA en estudios.}

\begin{titlebox}{milcMagenta}
Fase 3 · Praxis con pensamiento computacional
\end{titlebox}

\begin{softbox}{milcMagenta}{milcGris}{Construyendo el producto con los cuatro pilares}
\textbf{Actividad 3 · CREA --- Mi política personal de uso de IA en estudios} (28 min · individual). Vas a escribir el documento que regirá tu uso de IA en estudios durante el año académico.
\end{softbox}

\small
\begin{tabularx}{\textwidth}{>{\bfseries\RaggedRight\arraybackslash}p{3.6cm}Y}
\rowcolor{milcMagenta!90}\color{white}Pilar & \color{white}\bfseries Aplicación al producto\\
Descomposición & \textbf{Paso 1 --- Título y dedicatoria.} En el cuaderno escribe: \textbf{``Mi política personal de uso de IA en estudios''}. Debajo: \emph{``[Tu nombre], grado 7°, año académico 2026''}.\\
Patrones & \textbf{Paso 2 --- Los 4 usos LEGÍTIMOS con MI ejemplo concreto.} Lista numerada con: nombre del uso + ejemplo específico de cómo lo aplicaré. Ejemplo: \emph{``1. Como tutor: si no entiendo derivadas, le pido a ChatGPT que me las explique con un ejemplo del fútbol. Después escribo el resumen con mis palabras en el cuaderno. La IA ayuda a entender, no a entregar''}. Adapta los 4 a tu vida.\\
Abstracción & \textbf{Paso 3 --- Los 4 usos NO LEGÍTIMOS que me COMPROMETO a NO hacer.} Lista numerada con: nombre + por qué no lo haré. Ejemplo: \emph{``1. NO escribiré ensayos completos con IA y los pasaré como míos, porque eso me quita el oficio de aprender a escribir y faltaría a la honestidad académica''}. Adapta los 4 a tu vida.\\
Algoritmo de armado & \textbf{Paso 4 --- Plantilla de atribución que voy a usar.} Escribe en el cuaderno la plantilla concreta que vas a usar cuando uses IA en un trabajo. Ejemplo: \emph{``Asistencia de IA: en este trabajo usé [nombre del LLM] para [propósito específico]. Las ideas y el desarrollo son míos.''}. Guarda esa plantilla; la copiarás en futuros trabajos.\\
\end{tabularx}
\normalsize

\begin{drawbox}{milcMagenta}{Antes de cerrar la S9}
\checkbox Mi política tiene los 4 usos legítimos con ejemplos concretos míos. \\[1mm]
\checkbox Tengo los 4 usos no legítimos con compromiso de no hacerlos. \\[1mm]
\checkbox La plantilla de atribución está lista para copiar y pegar. \\[1mm]
\checkbox Las 5 preguntas de autoevaluación están subrayadas. \\[1mm]
\checkbox Hay ejemplo aplicado a un caso real o imaginado. \\[1mm]
\checkbox La firma comprometida está al final.
\end{drawbox}

\begin{softbox}{milcMagenta}{milcGris}{Plantilla de trabajo}
\textbf{Estructura.} \textit{Paso 1:} título + dedicatoria. \textit{Paso 2:} 4 legítimos con mi ejemplo. \textit{Paso 3:} 4 no legítimos con compromiso. \textit{Paso 4:} plantilla de atribución. \textit{Paso 5:} 5 preguntas subrayadas. \textit{Paso 6:} ejemplo aplicado + firma.
\end{softbox}

\begin{infoband}{milcMagenta}
\textbf{Lo que va al cuaderno (Actividad 3 · CREA).} Título: «Actividad 3 --- Mi política personal de IA en estudios». \textbf{Formato:} política completa de 6 secciones. \textbf{Extensión:} 1 página entera. \textbf{Sabes que terminaste cuando} sabes exactamente qué hacer y qué no hacer con IA en cada situación académica.
\end{infoband}

\puente{El producto es esa política firmada + ejemplo aplicado.}

\begin{titlebox}{milcMostaza}
Producto de la sesión
\end{titlebox}
\textbf{Política personal de uso de IA en estudios · 6 secciones · firmada comprometida}.
\begin{softbox}{milcMostaza}{milcGris}{Criterios de calidad}
(1) Los 4 usos legítimos tienen ejemplos concretos del estudiante. (2) Los 4 usos no legítimos tienen compromiso explícito. (3) La plantilla de atribución está lista para usar. (4) Las 5 preguntas de autoevaluación están subrayadas. (5) Hay ejemplo aplicado a un caso real o imaginado.
\end{softbox}

\puente{Antes de salir, intercambia tu política con un compañero --- ¿se parecen?.}

\begin{titlebox}{milcVino}
Fase 4 · Evaluación liberadora — cinco dimensiones
\end{titlebox}

\begin{softbox}{milcVino}{milcGris}{Sentido de la evaluación}
Evaluar no es solo revisar una nota. Es reconocer qué aprendí, qué cambió en mí, qué emociones manejé, cómo me relaciono con la ciudad y la comunidad, y qué le dejaré a quien viene detrás.
\end{softbox}

\textbf{Mi reflexión liberadora en cinco dimensiones.} Completa cada frase iniciada:

\small
\begin{tabularx}{\textwidth}{>{\bfseries\RaggedRight\arraybackslash}p{3.4cm}Y>{\RaggedRight\arraybackslash}p{4.6cm}}
\rowcolor{milcVino}\color{white}Dimensión & \color{white}\bfseries Mi respuesta (frame starter) & \color{white}\bfseries Algo que descubrí\\
1. Personal & Lo que más cambió en mí fue \linead{6.5cm} & \linead{4.2cm}\\
2. Emocional & La emoción más fuerte fue \linead{2.5cm} y la regulé \linead{3cm} & \linead{4.2cm}\\
3. Ciudadana & En democracia esto importa porque \linead{6cm} & \linead{4.2cm}\\
4. Local (Cartago) & En mi barrio sería útil para \linead{6.5cm} & \linead{4.2cm}\\
5. Intergeneracional & A mi abuela/abuelo le diría \linead{6.5cm} & \linead{4.2cm}\\
\end{tabularx}
\normalsize

\begin{infoband}{milcVino}
\textbf{Para la próxima sustentación.} Vuelve a esta tabla en seis meses. Verás que las respuestas cambian — no porque lo aprendido cambie, sino porque tú cambias. La evaluación liberadora es una conversación contigo a través del tiempo.
\end{infoband}


\puente{Cerramos con 3 voces sobre por qué el aprendizaje real no se reemplaza con IA.}

\begin{titlebox}{milcGradeDark}
Triángulo de pensamiento · cierre filosófico
\end{titlebox}

\begin{softbox}{milcVino}{milcGris}{Dussel · lente del nosotros}
\textit{````Lo que aprendes con esfuerzo te pertenece. Lo que obtienes con atajos, no es tuyo. El esfuerzo es lo que convierte el conocimiento en sabiduría propia.''''}

Pedro y Lucía de la historia de doña Mercedes pasaron de grado igual. Pero \emph{el conocimiento que Pedro tenía era suyo}; el de Lucía no. \textbf{Cuando usas IA para evitar esfuerzo, te haces Lucía}. Cuando la usas como asistente legítimo (tutor, corrector, organizador, generador de ideas), \emph{te haces Pedro con asistente}. La diferencia se acumula año tras año. \emph{Cada decisión sobre IA en tareas es voto por tu yo de 20 años: ¿quiero saber o quiero solo haber pasado de grado?}.

\textbf{Cuando uso IA en tareas, ¿estoy votando por convertirme en Pedro con asistente o en Lucía con atajos?}
\end{softbox}
\linead{\linewidth}\\[1mm]
\linead{\linewidth}

\begin{softbox}{milcMostaza}{milcGris}{Estoicismo · Marco Aurelio (emperador de la integridad) · cuidado interior}
\textit{````La integridad es la única posesión que no te pueden quitar. Pero tú mismo puedes regalarla con tus decisiones diarias.''''}

Cada vez que decides usar IA con criterio o sin él, \textbf{estás formando o desformando tu integridad personal}. Los profes no siempre detectan trampas; los compañeros pueden no saber. Pero \emph{tú sí sabes}. Esa conciencia interior es lo que el estoico cultiva. \emph{Tu política personal es promesa contigo mismo, no con el profe}. Cumplirla te hace persona íntegra; romperla te degrada.

\textbf{¿Estoy cultivando mi integridad o regalándola con decisiones pequeñas?}
\end{softbox}
\linead{\linewidth}\\[1mm]
\linead{\linewidth}

\begin{softbox}{milcTurquesa}{milcGris}{Floridi · lente de la infoesfera}
\textit{````Los estudiantes que aprendan a usar IA con criterio en grado 7 serán los profesionales del futuro. Los que solo copien, serán reemplazados por IAs mejores.''''}

\textbf{El mercado laboral del futuro no premia a quien copia de IA}: premia a quien usa IA como asistente y agrega valor humano genuino. Si solo entregas lo que la IA dice, \emph{cualquier persona con acceso a IA puede hacer lo mismo}. Tu valor profesional viene de \emph{tu criterio, tu voz, tu capacidad de evaluar y mejorar lo que la IA produce}. Aprender eso en 7° grado es invertir años de carrera profesional.

\textbf{¿Estoy formando habilidades que la IA no puede reemplazar, o me estoy haciendo reemplazable?}
\end{softbox}
\linead{\linewidth}\\[1mm]
\linead{\linewidth}

\begin{titlebox}{milcVino}
Compromiso final
\end{titlebox}

\small
\begin{tabularx}{\textwidth}{>{\RaggedRight\arraybackslash}p{3.0cm}YYY}
\rowcolor{milcVino}\color{white}\bfseries Criterio & \color{white}\bfseries Lo logré & \color{white}\bfseries En proceso & \color{white}\bfseries Necesito apoyo\\
Comprensión del tema & Explico ideas centrales con mis palabras. & Reconozco ideas pero falta precisión. & Necesito repasar conceptos base.\\
Pensamiento computacional & Apliqué los 4 pilares al diseñar mi producto. & Apliqué algunos pilares. & Necesito releer la fase 2.\\
Aplicación y producto & Mi producto cumple los criterios y se puede revisar. & Producto funcional con ajustes pendientes. & Aún en construcción.\\
Reflexión 5 dimensiones & Respondí cada dimensión con honestidad. & Respondí algunas con detalle. & Mis respuestas son muy generales.\\
Triángulo de pensamiento & Conecté las 3 voces con mi trabajo real. & Conecté al menos una. & No relaciono las voces con mi proceso.\\
\end{tabularx}
\normalsize

\textbf{Hoy entendí que:}\\[1mm]
\linead{\linewidth}\\[1mm]
\linead{\linewidth}

\textbf{Mi compromiso para la próxima sesión es:}\\[1mm]
\linead{\linewidth}\\[1mm]
\linead{\linewidth}

\begin{softbox}{milcMostaza}{milcGris}{Cita de despedida · Marco Aurelio}
\textit{``El obstáculo en el camino se vuelve el camino. Lo que estorba al camino se vuelve camino.''}\\[1mm]
Cada error de hoy, bien observado, se convierte en pieza del camino que pisarás mañana.
\end{softbox}

\small
\begin{tabularx}{\textwidth}{>{\bfseries}p{3.6cm}Y}
\rowcolor{milcGradeDark!12}Nombre completo: & \linead{10cm}\\
Fecha: & \linead{4cm} \quad Curso/grupo: \linead{4cm}\\
Mi firma: & \linead{9cm}\\
\end{tabularx}
\normalsize

\begin{infoband}{milcGradeDark}
\textbf{Para la próxima sesión.} Esta semana, en cada tarea donde uses IA, aplica tu política. La próxima clase es \textbf{la cosecha del periodo}: un proyecto completo con IA responsable. Vas a aplicar todo lo aprendido en P3 a un trabajo real con atribución, reflexión ética y sustentación.
\end{infoband}

\end{document}
